\subsection{极小投射分解}

设 M 是一个 A-模，且设

\[
\dots \longrightarrow \mathrm{P} _ {n} \xrightarrow {d _ {n}} \mathrm{P} _ {n - 1} \longrightarrow \dots \longrightarrow \mathrm{P} _ {0} \xrightarrow {d _ {0}} \mathrm{M} \longrightarrow 0\tag{P}
\]

为M的一个分解。我们称(P)为极小投射分解，如果对每一个 \(n \geqslant 0\)，同态 \(\delta_{n}: P_{n} \to \operatorname{Im}(d_{n})\)（由 \(d_{n}\) 所诱导）是一个投射覆盖（VIII，§8，第5号）。

命题8.—设M是一个A-模，且P与P'是M的两个极小投射分解，其中f：P→P'是分解的一个态射。则f是一个同构。特别地，M的任意两个极小投射分解都是同构的。

令 \(\tilde{\mathbf{P}}_{n} = \mathbf{P}_{n}\) 于 \(n \neq -1\)，且 \(\tilde{\mathbf{P}}_{-1} = \mathbf{M}\)；让我们类似地定义 \(\tilde{\mathbf{P}}_{n}^{\prime}\) 并设 \(f_{-1} = 1_{\mathbf{M}}\)。让我们从 \(-1\) 出发用归纳法证明 \(f_{n}: \tilde{\mathbf{P}}_{n} \to \tilde{\mathbf{P}}_{n}^{\prime}\) 对所有 \(n\) 都是同构。这对 \(n = -1\) 是显然的；假设 \(f_{n}\) 与 \(f_{n-1}\) 是同构。由图的交换性可得：

\[
\begin{array}{c c c} \mathbf {P} _ {n} & \xrightarrow {d _ {n}} & \mathbf {P} _ {n - 1} \\ \Big \downarrow f _ {n} & & \Big \downarrow f _ {n - 1} \\ \mathbf {P} _ {n} ^ {\prime} & \xrightarrow {d _ {n} ^ {\prime}} & \mathbf {P} _ {n - 1} ^ {\prime} \end{array}
\]

即 \(f_{n}\) 诱导出 \(g_{n}\)，它是从 \(\operatorname{Ker} d_{n}\) 到 \(\operatorname{Ker} d_{n}^{\prime}\) 上的一个同构。于是由图的交换性以及

\[
\begin{array}{c c c} \mathrm{P} _ {n + 1} & \xrightarrow {\delta_ {n + 1}} & \text {Ker} d _ {n} \\ \Big \downarrow f _ {n + 1} & & \Big \downarrow g _ {n} \\ \mathrm{P} _ {n + 1} ^ {\prime} & \xrightarrow {\delta_ {n + 1} ^ {\prime}} & \text {Ker} d _ {n} ^ {\prime} \end{array}
\]

以及 VIII，前引文献，可得 \(f_{n+1}\) 是一个同构。

推论. --- 设 M 是一个 A-模，P 和 P' 是 M 的两个投射分解；假设 P 是极小的。设 f : P → P' 和 g : P' → P 是两个分解态射。则 f 是单射，g 是满射，且 P' 是子复形 Im f 与 Ker g 的直和。此外，Ker g 具有零同调。

事实上， \(\alpha = g \circ f\) 是P的一个自同构（命题8）。设 \(\tilde{f} = f \circ \alpha^{-1}\) 。我们有

\[
\operatorname{Im} \tilde {f} = \operatorname{Im} f \quad \text { and } \quad g \circ \tilde {f} = 1 _ {\mathrm{P}},
\]

这表明 \(\mathbf{P}' = \operatorname{Im} \tilde{f} \oplus \operatorname{Ker} g\) 。由于序列

\[
0 \rightarrow \operatorname{Ker} g \rightarrow \mathrm{P} ^ {\prime} \stackrel {g} {\rightarrow} \mathrm{P} \rightarrow 0
\]

是正合的，且 \(g\) 是一个拟同构， \(\operatorname{Ker} g\) 具有零同调。

命题9。——设M是一个A-模，且(P, p)是M的一个投射分解。设r表示A的根。假设以下任一条件成立： \(P_{n}\) 对所有n都是有限生成的A-模，或者r是幂零的。那么 (P, p) 为极小的充要条件是复形 \((\mathbf{A}/\mathbf{r}) \otimes_{\mathbf{A}} \mathbf{P}\) 具有零微分，换言之即

\[
d _ {n + 1} (\mathbf {P} _ {n + 1}) \subset \mathrm{r} \mathbf {P} _ {n} \quad \text { for all } n \geqslant 0.
\]

假设 (P, p) 是极小的。根据 VIII，前引文献，同态

\[
1 \otimes \delta_ {n}: (\mathrm{A} / \mathrm{r}) \otimes_ {\mathrm{A}} \mathrm{P} _ {n} \rightarrow (\mathrm{A} / \mathrm{r}) \otimes_ {\mathrm{A}} \operatorname{Im} d _ {n}
\]

是一个同构。于是由正合序列可得

\[
0 \longrightarrow \operatorname{Im} d _ {n + 1} \xrightarrow {j _ {n}} \mathrm{P} _ {n} \xrightarrow {\delta_ {n}} \operatorname{Im} d _ {n} \longrightarrow 0
\]

使得同态 \(1 \otimes j_n : (\mathbf{A} / \mathfrak{r}) \otimes_{\mathbf{A}} \operatorname{Im} d_{n+1} \to (\mathbf{A} / \mathfrak{r}) \otimes_{\mathbf{A}} \mathrm{P}_n\) 为零；由于 \(d_{n+1} = j_n \circ \delta_{n+1}\) ，我们推出 \(1_{\mathbf{A} / \mathfrak{r}} \otimes d_{n+1} = 0\) ，对 \(n \geqslant 0\) 。

反之，假设对所有 \(n \geqslant 1\) , \(1 \otimes d_n\) 为零，即 \(\operatorname{Im} d_n = \operatorname{Ker} d_{n-1}\) 包含于 \(\mathbf{rP}_{n-1}\) 。由于 \(\delta_{n-1}\) 是满射的，由VIII，前引文献可知 \(\delta_{n-1}\) 对 \(n \geqslant 1\) 是一个投射覆盖，因此 \((\mathbf{P}, p)\) 是极小的。

命题10。——设A是左诺特局部环，令M为有限生成的A模。则M具有极小分解(P, p)；对每个 \(n \geqslant 0\) , \(\mathbf{P}_{n}\) 是有限生成的自由模。

事实上，在第5节（第53页）所作的构造中，由VIII，前引文献，可以取 \((\mathrm{L}_{0}, p)\) 为 M 的一个投射覆盖，并取 \(L_{n+1}\) 为 Ker \(d_{n}\) 的一个投射覆盖。如此得到的分解便是极小的。

注记. --- 1) 设 m 表示 A 的极大理想，并设 \(k = \mathbf{A}/\mathfrak{m}\) 。设 P 为 M 的一个极小投射分解，并设 \(b_{n} = \dim_{k}(k \otimes_{\mathbf{A}} \mathbf{P}_{n})\) 。则 \(\mathbf{P}_{n}\) 是秩为 \(b_{n}\) 的自由 A-模。由命题 8 的推论可知，对于 M 的任何其他投射分解 \(\mathbf{P}'\) ，有 \(\dim_{k}(k \otimes_{\mathbf{A}} \mathbf{P}_{n}') \geqslant b_{n}\) ，且等号成立当且仅当 \(\mathbf{P}'\) 是极小的。

\begin{enumerate}
\def\labelenumi{\arabic{enumi})}
\setcounter{enumi}{1}
\tightlist
\item
  由命题 9， \(b_n\) 是 \(k\) 上 \(H_n(k \otimes_A P)\) 的维数，* 换言之即 \(\text{Tor}_n^A(k, M)\) 的维数。它也是 \(k\) 上 \(\text{Ext}_A^n(M, k)\) 的维数（参见 X，第 103 页，注记 3）*。
\end{enumerate}

